% ============================================================
%  导言区：样式与宏包（CUMCM 论文通用模板，实测可编译）
%  依据：官方《全国大学生数学建模竞赛论文格式规范》
%        —— 页边距上下左右 >=2.5cm；字体/字号/行距官方未强制
% ============================================================

% ---------- 页面 ----------
\usepackage{geometry}
\geometry{a4paper, top=2.5cm, bottom=2.5cm, left=2.5cm, right=2.5cm}

% ---------- 行距（为满足正文 <=30 页限制，由 1.5 调整为 1.38） ----------
\linespread{1.38}

% ---------- 页码：页脚居中（官方规范） ----------
\pagestyle{plain}

% ---------- 章节编号：中文数字 + 、（如 “一、问题重述”） ----------
\ctexset{
  section = {number = \chinese{section}, aftername = {、}},
}

% ---------- 字体：正文中文宋体（ctex 默认）；西文/数字/公式为 Times（MathType 观感） ----------
\setmainfont{Times New Roman}

% ---------- 数学 ----------
\usepackage{amsmath, amsfonts, mathtools}
\usepackage{unicode-math}
\setmathfont{XITSMath-Regular.otf}   % XITS Math = Times 系 OpenType 数学字体

% ---------- 定理环境 ----------
\usepackage{amsthm}
\newtheorem{lemma}{引理}
\newtheorem{theorem}{定理}
\newtheorem{proposition}{命题}
\newtheorem{corollary}{推论}

% ---------- 图表 ----------
\usepackage{graphicx}
\graphicspath{{figures/}}
\usepackage{booktabs}
\usepackage{multirow}
\usepackage{makecell}
\usepackage{threeparttable}
\usepackage{tabularx}
\usepackage{longtable}
\usepackage{float}
\usepackage{caption}
\usepackage{subcaption}
\captionsetup{font=small, labelsep=quad}

% ---------- 列表 ----------
\usepackage{enumitem}

% ---------- 颜色与代码 ----------
\usepackage{xcolor}
\usepackage{listings}
\lstset{
  basicstyle=\ttfamily\small,
  breaklines=true,
  frame=single,
  numbers=left,
  numberstyle=\tiny,
}

% ---------- 算法伪代码 ----------
\usepackage[ruled, vlined, linesnumbered]{algorithm2e}
\SetAlgorithmName{算法}{算法}{算法索引}

% ---------- 单位 ----------
\usepackage{siunitx}

% ---------- 参考文献（GB/T 7714—2015 顺序编码制） ----------
\usepackage{natbib}
\usepackage{gbt7714}

% ---------- 超链接（放在最后加载） ----------
\usepackage[colorlinks=true, linkcolor=black, citecolor=black, urlcolor=blue]{hyperref}
\usepackage{xurl}

% ---------- 浮动体版面参数（减少整页空白） ----------
\renewcommand{\topfraction}{0.9}
\renewcommand{\bottomfraction}{0.85}
\renewcommand{\textfraction}{0.07}
\renewcommand{\floatpagefraction}{0.72}
\setcounter{topnumber}{3}
\setcounter{bottomnumber}{2}
\setcounter{totalnumber}{4}
